\subsection{The Challenge of Manual Memory Management}

In the early days of programming, memory management was entirely manual. Languages such as C and C++ required programmers to explicitly control memory allocation and deallocation, providing precise control over system resources at the cost of increased development complexity. This manual approach requires programmers to allocate memory explicitly (using functions like \texttt{malloc} in C) and later free that memory when it is no longer needed (using \texttt{free}).

While manual memory management offers fine-grained control and potentially optimal performance, it introduces significant risks:

\begin{itemize}
    \item \textbf{Memory leaks} occur when allocated memory is never freed, gradually depleting available system resources.
    \item \textbf{Dangling pointers} arise when a program accesses memory that has already been freed, causing unpredictable behavior.
    \item \textbf{Double frees} happen when memory is deallocated multiple times, potentially corrupting memory management data structures.
    \item \textbf{Buffer overflows} result from writing beyond allocated memory boundaries, creating security vulnerabilities.
\end{itemize}

The prevalence of these errors in production software has driven the evolution toward safer memory management alternatives.

\subsubsection*{Manual Memory Management Examples in C}

\begin{lstlisting}[language=C, caption=Memory Leak Example, style=mystyle]
#include <stdlib.h>

void memory_leak() {
    int* data = (int*) malloc(100 * sizeof(int));
    // Function ends without freeing 'data'
    // This memory remains allocated but inaccessible
}
\end{lstlisting}

\begin{lstlisting}[language=C, caption=Dangling Pointer Example, style=mystyle]
#include <stdlib.h>
#include <stdio.h>

void dangling_pointer() {
    int* data = (int*) malloc(sizeof(int));
    *data = 42;
    free(data);
    // 'data' now points to invalid memory
    printf("%d\n", *data);  // Undefined behavior!
}
\end{lstlisting}

\subsection{The Evolution to Automatic Memory Management}

As software systems grew in complexity, the burden of manual memory management became increasingly problematic. This led to the development of automatic memory management techniques that could reliably handle memory allocation and deallocation without programmer intervention. Two primary approaches emerged: reference counting and tracing garbage collection.

\subsubsection{Reference Counting: A First Step}

Reference counting represents an intermediate step between manual and fully automatic memory management. In this approach, each allocated object maintains a counter that tracks how many references point to it:

\begin{itemize}
    \item When a new reference to an object is created, its reference count increments.
    \item When a reference is removed, the count decrements.
    \item When the count reaches zero, the object is automatically deallocated.
\end{itemize}

\begin{lstlisting}[language=Python, caption=Reference Counting Concept, style=pythonstyle]
def create_reference(obj):
    obj.ref_count += 1

def remove_reference(obj):
    obj.ref_count -= 1
    if obj.ref_count == 0:
        # Object has no more references
        for child in obj.references:
            remove_reference(child)  # Decrement children
        free(obj)  # Deallocate the object
\end{lstlisting}

While reference counting immediately reclaims memory when objects become unreachable and distributes collection overhead throughout program execution, it has a critical limitation: it cannot detect cyclic references.

\begin{figure}[h]
\centering
\begin{tikzpicture}[
  every node/.style={circle, draw, minimum size=1.2cm},
  node distance=2.5cm,
  arrow/.style={->, thick}
]
\node (A) {A: 1};
\node (B) [right=of A] {B: 1};

\draw[arrow] (A) -- (B) node[midway, above, draw=none] {};
\draw[arrow] (B) to [bend left=30] (A) node[midway, below, draw=none] {};

\node[draw=none] at (0,-2) {Even though neither A nor B is reachable from};
\node[draw=none] at (0,-2.5) {program roots, their reference counts remain at 1.};
\end{tikzpicture}
\caption{The cyclic reference problem: Objects A and B reference each other but are unreachable from program roots. Reference counting alone cannot detect this situation.}
\label{fig:cyclic_refs}
\end{figure}

As shown in Figure \ref{fig:cyclic_refs}, when objects form reference cycles, they may maintain non-zero reference counts even when the entire cycle is unreachable from the rest of the program. This results in memory leaks, undermining one of the primary benefits of automatic memory management.

\subsection{Tracing Garbage Collection: A Comprehensive Solution}

To address the limitations of reference counting, tracing garbage collection was developed based on a fundamentally different principle: \textbf{reachability analysis}. Rather than tracking individual references, tracing collectors determine whether objects are reachable from a set of root references within the program.

\subsubsection{Reachability Analysis: The Foundation of Modern GC}

In reachability analysis, memory can be modeled as a directed graph where:
\begin{itemize}
    \item Nodes represent allocated objects
    \item Edges represent references between objects
    \item Special nodes called "roots" represent references directly accessible to the program (e.g., global variables, stack variables, and registers)
\end{itemize}

The key insight of reachability analysis is that an object is "live" (still needed) if and only if it is reachable by following a path of references from any root. Conversely, any object that cannot be reached from the roots is definitionally unreachable by the program and can safely be reclaimed.

\begin{figure}[h]
\centering
\begin{tikzpicture}[
  every node/.style={circle, draw, minimum size=1.2cm},
  node distance=2.2cm and 2.5cm,
  arrow/.style={->, thick},
  reachable/.style={fill=gray!20},
  unreachable/.style={fill=red!15}
]

\node[reachable] (root) {Root};
\node[reachable] (A) [right=of root] {A};
\node[reachable] (B) [below=of A] {B};
\node[reachable] (C) [right=of A] {C};
\node[reachable] (D) [below=of C] {D};
\node[unreachable] (E) [right=of D] {E};
\node[unreachable] (F) [right=of C] {F};

\draw[arrow] (root) -- (A);
\draw[arrow] (A) -- (B);
\draw[arrow] (A) -- (C);
\draw[arrow] (C) -- (D);
\draw[arrow] (E) -- (F);
\draw[arrow] (F) -- (E);

\end{tikzpicture}
\caption{Reachability analysis in garbage collection. Objects E and F form a cycle but are unreachable from any root, making them eligible for collection despite their mutual references.}
\label{fig:reachability}
\end{figure}

As illustrated in Figure \ref{fig:reachability}, reachability analysis elegantly solves the cyclic reference problem. Even though objects E and F reference each other (maintaining non-zero reference counts), they are unreachable from any root and therefore eligible for collection.

\subsubsection{Mark-and-Sweep: The Archetypal Tracing Algorithm}

The most fundamental tracing garbage collection algorithm is Mark-and-Sweep, which operates in two distinct phases:

\begin{itemize}
    \item \textbf{Mark Phase}: Starting from all roots, traverse the object graph and mark every object encountered as "reachable."
    \item \textbf{Sweep Phase}: Scan the entire heap and reclaim any unmarked objects, as they cannot be reached from the roots.
\end{itemize}

\begin{lstlisting}[language=Python, caption=Mark-and-Sweep Algorithm, style=pythonstyle]
def mark_and_sweep():
    # Mark phase
    for root in program_roots:
        mark(root)
    
    # Sweep phase
    for object in heap:
        if not object.is_marked:
            free(object)  # Object is unreachable, reclaim it
        else:
            object.is_marked = False  # Reset for next collection

def mark(object):
    if object is None or object.is_marked:
        return  # Already marked or null reference
    
    object.is_marked = True  # Mark as reachable
    
    # Recursively mark all referenced objects
    for reference in object.references:
        mark(reference)
\end{lstlisting}

The mark-and-sweep approach provides several advantages:
\begin{itemize}
    \item It correctly identifies and collects cyclic garbage
    \item It guarantees that all unreachable objects will be reclaimed
    \item It can be implemented efficiently with various optimizations
\end{itemize}